
\documentclass[11pt]{article}
\input{../shared/project_style.tex}
\rhead{Module 5}
\begin{document}

\DocTitle{Module 5: Future Nonequilibrium Electron-Phonon and Photon Physics}{Photon absorption, carrier excitation, relaxation, recombination, and pair breaking}{Physics note for FEM\_BallisticDiffusion\_PINN}

\begin{overviewbox}
\textbf{Purpose.} This note documents physics that should be added after the first FEM and PINN framework is stable. No code is required in this module for the first build. The goal is to create a clear roadmap for future nonequilibrium physics: photon absorption, hot-electron generation, electron-phonon relaxation, high-energy phonon creation, Cooper-pair breaking, quasiparticle recombination, and phonon bottlenecks.
\end{overviewbox}

\ProjectTOC

\SymbolsAcronymsSection
\begin{symtable}
$h\nu$ & Photon energy & Energy of an absorbed photon of frequency $\nu$. \\
$\Delta$ & Superconducting energy gap & Minimum quasiparticle excitation energy scale. \\
$2\Delta$ & Pair-breaking threshold & Minimum energy required to break a Cooper pair into two quasiparticles. \\
$N_{\qp}$ & Quasiparticle number & Total number of quasiparticles in a modeled region. \\
$N_\Omega$ & Pair-breaking phonon number & Number of high-energy phonons capable of breaking Cooper pairs. \\
$\tau_R$ & Recombination time & Characteristic time for quasiparticle recombination. \\
$\tau_B$ & Pair-breaking time & Characteristic time for high-energy phonons to break pairs. \\
$\tau_{\mathrm{esc}}$ & Phonon escape time & Time for high-energy phonons to leave the superconducting film or escape into substrate/sink. \\
$\Sigma$ & Electron-phonon coupling coefficient & Material parameter in reduced electron-phonon heat-transfer laws. \\
$T_e$ & Electron temperature & Effective temperature of electron/quasiparticle subsystem. \\
$T_{ph}$ & Phonon temperature & Effective temperature of phonon/lattice subsystem. \\
RT & Rothwarf-Taylor & Coupled rate-equation model for quasiparticles and pair-breaking phonons. \\
\end{symtable}

\section{Why this is future work}
Modules 1--4 create a usable reduced transport framework. Module 5 is more microscopic and more nonequilibrium. It involves energy-resolved populations, pair breaking, phonon escape, and electron-phonon relaxation. These processes are essential for a mature model, but adding them before the geometry/FEM framework works would make the project too complicated too early.

The future extension should answer questions such as:
\begin{enumerate}[leftmargin=1.6em]
    \item Where is a photon absorbed?
    \item How much of the photon energy creates quasiparticles?
    \item How much becomes high-energy phonons?
    \item Do those phonons escape, thermalize, or break more Cooper pairs?
    \item How do traps and heat sinks alter these pathways?
\end{enumerate}

\section{Photon absorption and pair-breaking threshold}
A photon with energy
\begin{equation}
E_\gamma=h\nu
\end{equation}
can break a Cooper pair if it can supply at least
\begin{equation}
E_\gamma\ge 2\Delta.
\end{equation}
The factor of two appears because a Cooper pair must be broken into two quasiparticle excitations, each costing at least roughly $\Delta$ in energy.

Not every absorbed photon produces the same number of quasiparticles. Some energy may remain in hot electrons, some may go into high-energy phonons, and some may escape before generating additional quasiparticles. A reduced source term can be written as
\begin{equation}
S_{\qp}(\rvec,t)=\eta_{\qp}\frac{Q_{\mathrm{abs}}(\rvec,t)}{\Delta},
\end{equation}
where $Q_{\mathrm{abs}}$ is absorbed power density and $\eta_{\qp}$ is an efficiency factor. The factor $\Delta$ converts energy rate density into approximate quasiparticle generation rate density.

\section{Electron-phonon relaxation}
After energy is deposited in a metal film, the electronic subsystem and phonon subsystem may not immediately share one temperature. A reduced two-temperature model is
\begin{equation}
C_e(T_e)\frac{\partial T_e}{\partial t}=\nabla\cdot(k_e\nabla T_e)-G_{e-ph}(T_e,T_{ph})+Q_{\mathrm{abs}},
\end{equation}
\begin{equation}
C_{ph}(T_{ph})\frac{\partial T_{ph}}{\partial t}=\nabla\cdot(k_{ph}\nabla T_{ph})+G_{e-ph}(T_e,T_{ph})-G_{\mathrm{esc}}(T_{ph},T_{\mathrm{sub}}).
\end{equation}
At low temperature, a common normal-metal electron-phonon energy-transfer form is
\begin{equation}
G_{e-ph}=\Sigma\left(T_e^5-T_{ph}^5\right).
\end{equation}
In a superconductor, the gapped density of states modifies these laws. Therefore this form should be treated as a conceptual bridge, not a final superconducting microscopic model.

\section{Rothwarf-Taylor rate equations}
A classic reduced model for nonequilibrium quasiparticles and high-energy phonons is the Rothwarf-Taylor system. In a spatially lumped form,
\begin{equation}
\frac{\dd N_{\qp}}{\dd t}=-R N_{\qp}^2+2\beta N_\Omega+G_{\qp},
\end{equation}
\begin{equation}
\frac{\dd N_\Omega}{\dd t}=\frac{1}{2}R N_{\qp}^2-\beta N_\Omega-\frac{N_\Omega}{\tau_{\mathrm{esc}}}+G_\Omega.
\end{equation}
Here $N_\Omega$ is the number of high-energy phonons capable of breaking Cooper pairs, $R$ is a recombination coefficient, $\beta$ is a pair-breaking rate by high-energy phonons, $\tau_{\mathrm{esc}}$ is the phonon escape time, and $G_{\qp}$ and $G_\Omega$ are generation terms.

The physical feedback loop is important. Recombination removes two quasiparticles, but it emits a phonon with energy near $2\Delta$. If that phonon remains in the film, it can break another Cooper pair. This produces a phonon bottleneck. A good heat sink or substrate escape pathway can reduce this bottleneck by removing high-energy phonons.

\section{Spatial extension}
A spatially resolved future model could couple quasiparticle diffusion and high-energy phonon diffusion:
\begin{equation}
\frac{\partial n_{\qp}}{\partial t}
=\nabla\cdot(D_{\qp}\nabla n_{\qp})-R n_{\qp}^2+2\beta n_\Omega+S_{\qp}-\Gamma_{\mathrm{trap}}n_{\qp},
\end{equation}
\begin{equation}
\frac{\partial n_\Omega}{\partial t}
=\nabla\cdot(D_\Omega\nabla n_\Omega)+\frac{1}{2}R n_{\qp}^2-\beta n_\Omega-\frac{n_\Omega}{\tau_{\mathrm{esc}}}+S_\Omega.
\end{equation}
This is a natural Module 5 future upgrade because it connects directly to the geometry already built in Module 3.

\section{Relaxation and recombination cascade}
A realistic absorbed-energy event may pass through the following stages:
\begin{enumerate}[leftmargin=1.6em]
    \item photon or energetic phonon deposits energy in a film, substrate, or junction region,
    \item high-energy electronic excitations are created,
    \item electron-electron scattering redistributes energy inside the electronic subsystem,
    \item electron-phonon scattering transfers energy to phonons,
    \item phonons with energy above $2\Delta$ can break Cooper pairs,
    \item quasiparticles diffuse and recombine,
    \item recombination phonons either escape, down-convert, or break pairs again,
    \item traps and heat sinks modify the escape and removal rates.
\end{enumerate}

\section{How Module 5 connects back to earlier modules}
The future nonequilibrium model should not replace Modules 1--4. It should provide better source and coupling terms:
\begin{equation}
S_{\qp}(\rvec,t),\qquad Q_{\mathrm{abs}}(\rvec,t),\qquad \Gamma_{\mathrm{trap}}(\rvec),\qquad \tau_{\mathrm{esc}}(\rvec),\qquad \eta_{\qp}(\rvec).
\end{equation}
These improved terms can then be inserted into the FEM and PINN framework already established.

\section{First future-work milestones}
The recommended order for future work is:
\begin{enumerate}[leftmargin=1.6em]
    \item add a spatially localized photon absorption source term,
    \item add an energy-to-quasiparticle conversion efficiency,
    \item add a lumped Rothwarf-Taylor model coupled to the junction region,
    \item extend the lumped model to spatially resolved $n_{\qp}$ and $n_\Omega$,
    \item couple phonon escape to substrate and heat-sink geometry,
    \item train a PINN surrogate for the expanded nonequilibrium model.
\end{enumerate}

\section{References}
\begin{enumerate}[label={[\arabic*]}]
\item A. Rothwarf and B. N. Taylor, ``Measurement of recombination lifetimes in superconductors,'' \emph{Physical Review Letters} \textbf{19}, 27--30 (1967).
\item S. B. Kaplan et al., ``Quasiparticle and phonon lifetimes in superconductors,'' \emph{Physical Review B} \textbf{14}, 4854--4873 (1976).
\item M. Tinkham, \emph{Introduction to Superconductivity}, 2nd ed., Dover (2004).
\item P. J. de Visser et al., ``Number fluctuations of sparse quasiparticles in a superconductor,'' \emph{Physical Review Letters} \textbf{106}, 167004 (2011).
\end{enumerate}

\end{document}
